\documentclass[11pt]{article}
\usepackage[margin=0.8in]{geometry}
\usepackage{graphicx}
\usepackage{xcolor}
\usepackage{hyperref}
\usepackage{listings}

\hypersetup{
  colorlinks=true,
  linkcolor=blue,
  urlcolor=blue
}

\lstdefinestyle{pythonstyle}{
  language=Python,
  basicstyle=\ttfamily\small,
  breaklines=true,
  frame=single,
  columns=fullflexible,
  keepspaces=true,
  showstringspaces=false,
  keywordstyle=\color{blue!70!black},
  commentstyle=\color{green!40!black},
  stringstyle=\color{orange!70!black}
}

\title{Addition Self-Improvement Round 8 Model}
\author{}
\date{}

\begin{document}
\maketitle

\section*{What This Is}

This folder contains the final round-8 checkpoint from the fixed-width mixed-prompt
addition self-improvement run. The selected condition is filtered carry-aware
composition with a fixed binary composition path.

\begin{itemize}
  \item Architecture: recipe no-position LLaMA-style causal LM, 6 layers, 384 hidden units, 6 heads.
  \item Training regime: seed widths 3--7, then 8 self-improvement rounds.
  \item Final horizon: internal operand width 23.
  \item Final expanded held-out accuracy: 0.8644.
  \item Final frontier-train accuracy: 0.9653.
  \item Final seed-range held-out accuracy: 1.0000.
\end{itemize}

\section*{Input And Output Format}

The model expects a natural, unpadded arithmetic prompt:

\begin{quote}
\texttt{Q: \{a\} + \{b\} = ?}\\
\texttt{A:}
\end{quote}

For example:

\begin{quote}
\texttt{Q: 104 + 2305 = ?}\\
\texttt{A:}
\end{quote}

The output is an unpadded decimal integer answer. During training, examples were
sampled by an internal fixed operand width, but the visible prompt was always natural:
internal leading zeros were not shown to the model. Thus an internal width-2 example
such as \texttt{01 + 23} was displayed as \texttt{1 + 23}, with target \texttt{24}.

\section*{Loading And Inference}

This checkpoint uses the repository's custom recipe model class and arithmetic
character tokenizer. Load it through the local repository code rather than plain
vanilla \texttt{AutoModelForCausalLM}.

\begin{lstlisting}[style=pythonstyle]
from pathlib import Path
import torch

from self.core.recipes import (
    RECIPE_ARITHMETIC_SELF_IMPROVE_V1,
    build_recipe_tokenizer,
    load_recipe_model,
    resolve_addition_recipe,
    tokenizer_padding_side,
)
from self.core.evaluation import build_generation_encodings
from core.addition_pipeline import extract_numeric_answer

model_dir = Path("models/addition_self_improvement_round8")

preset = resolve_addition_recipe(RECIPE_ARITHMETIC_SELF_IMPROVE_V1)
tokenizer = build_recipe_tokenizer(preset)
model = load_recipe_model(
    model_dir,
    tokenizer,
    bf16=torch.cuda.is_available(),
    fp16=False,
)
model.eval()

def predict_addition(a: int, b: int, max_new_tokens: int = 48):
    prompt = f"Q: {int(a)} + {int(b)} = ?\nA:"
    device = next(model.parameters()).device
    with tokenizer_padding_side(tokenizer, "left"), torch.no_grad():
        enc = build_generation_encodings(tokenizer, [prompt], device)
        out = model.generate(
            **enc,
            max_new_tokens=max_new_tokens,
            do_sample=False,
        )
    generated = out[0, enc["input_ids"].shape[1]:]
    decoded = tokenizer.decode(generated, skip_special_tokens=True)
    return extract_numeric_answer(decoded), decoded

answer, raw_text = predict_addition(104, 2305)
print(answer, raw_text)
\end{lstlisting}

\section*{Held-Out Accuracy Heatmap}

Rows are internal operand widths and columns are self-improvement rounds. The visible
prompts remain natural/unpadded for every width.

\begin{center}
\includegraphics[width=0.95\linewidth]{heatmap_with_carry_filtered_round8.pdf}
\end{center}

\section*{Included Files}

\begin{itemize}
  \item \texttt{model.safetensors}: model weights.
  \item \texttt{config.json}: model configuration.
  \item \texttt{generation\_config.json}: generation configuration.
  \item \texttt{tokenizer\_config.json}: arithmetic character tokenizer metadata.
  \item \texttt{metrics.json}: final round metrics.
  \item \texttt{heatmap\_with\_carry\_filtered\_round8.pdf}: accuracy heatmap shown above.
\end{itemize}

\end{document}
